% Source: https://artofproblemsolving.com/community/c3h1379010_chinese_remainder_theorem
% Stable identity: aops-topic-1379010-q01
(a) $b_i$ is the product of the $m_j$ with $j\ne i$, each relatively prime to $m_i$, so $\gcd(b_i,m_i)=1$ and $b_i$ has an inverse $y_i$ modulo $m_i$. Fix $i$. For $j\ne i$, $b_j$ contains the factor $m_i$, so $a_jb_jy_j\equiv0\pmod{m_i}$. Hence
\[
x\equiv a_ib_iy_i\equiv a_i\cdot1=a_i\pmod{m_i}.
\]

(b) If $x'\equiv x\pmod M$, then $x'\equiv x\pmod{m_i}$ for each $i$ (as $m_i\mid M$), so $x'$ is a solution. Conversely, if $x'$ is a solution, then $m_i\mid x'-x$ for every $i$; the $m_i$ are pairwise relatively prime, so their product $M$ divides $x'-x$.

(c) $M=105$ and $b_1=35$, $b_2=21$, $b_3=15$. Since $35\equiv2\pmod3$, $y_1=2$; since $21\equiv1\pmod5$, $y_2=1$; since $15\equiv1\pmod7$, $y_3=1$. So
\[
x=2\cdot35\cdot2+3\cdot21\cdot1+2\cdot15\cdot1=140+63+30=233\equiv23\pmod{105}.
\]
The solutions are $x\equiv\boxed{23}\pmod{105}$. Check: $23=3\cdot7+2=5\cdot4+3=7\cdot3+2$.
